\section{Introduction}\label{sec:intro}

The Newman deformation combines two distinct mathematical objects: a
positive source on a real spectral line and an oscillatory entire
transform whose real zeros can move. An imaginary-time realization makes
the evolution of the source and the field precise, but its positivity
does not by itself control the nodal geometry of the observed amplitude.
To understand what additional information could be useful, it is helpful
first to describe a collision in the unchanged realization.

We use the normalization
\begin{equation}\label{eq:intro-H}
 H_\lambda(z)=\int_0^\infty \ee^{\lambda u^2}\Phi(u)\cos(zu)\dd u,
 \qquad H_0(z)=\frac18\xi\!\left(\frac{1+\ii z}{2}\right),
\end{equation}
where \(\xi(s)=s(s-1)\pi^{-s/2}\Gamma(s/2)\zeta(s)/2\).
The time convention is
\begin{equation}\label{eq:clock}
 h(\tau,z)=H_{\lambda_0-\tau}(z),\qquad
 \partial_\tau h=\partial_z^2h,\qquad \tau\geq0.
\end{equation}
Thus increasing imaginary time means decreasing Newman parameter.
The normalization and sign agree with the canonical source convention of
Rodgers and Tao \cite{RodgersTao2020}.

The analysis has three parts. First, we establish the complete source bounds
and realize \eqref{eq:clock} as a fixed observation of a contraction
semigroup on the full spherical-cylinder Hilbert space. The relative
generator is \(K=L_0-4I\). The ambient operator retains every angular shell
and both Dirac signs; the preparation occupies one declared reducing
channel. The source and its calibration are supplied, rather than
generated intrinsically by the geometry.

Second, we investigate a hypothetical event
\(h(\tau_*,x_*)=h_x(\tau_*,x_*)=0\). For a collision of two zeros, set
\(a=h_{xx}\ne0\), \(b=h_{xxx}\), \(c=h_{xxxx}\), all at the event.
Three descriptions then fit together:
\begin{itemize}
\item the individual roots have a square-root splitting, while the
critical value crosses zero smoothly and the squared gap is analytic;
\item the opposite-sign region between the roots loses both mass and
kinetic form, with different exponents, so its normalized local quotient
diverges although the global state is regular;
\item the collision is a pair of signed interference balances, with a
nonzero second-moment imbalance governing its temporal crossing.
\end{itemize}
The resulting formulas are local and event-dependent. Higher finite
multiplicity has a different leading heat polynomial, and is treated
separately. No assumption that theta collisions always have multiplicity two is made.

The basic local heat geometry is classical \cite{Angenent1991}. The
repository's earlier collision and multiplicity analyses supply the
immediate context \cite{TNDDirectCollision,HeatMultiplicity}. Here these
laws are developed together with the actual cylinder amplitude,
operator-domain bounds and local mass/energy asymptotics. The local
factorization also separates a pair's mutual motion from the regular
factor representing the surrounding field, without an infinite sum over
zeros. The distinction between field energy and singular zero-gap
quantities is essential.

Third, we ask whether the statistical contradiction can be pursued
when the zero configuration is not real. The answer begins with an
existing unconditional extension of Montgomery's theorem, rather
than with a new conjectural definition. We distinguish that complex
sum from ordinate-only and marked planar statistics, prove finite
exception sensitivity bounds, and connect the transverse motion of
a finite cluster to the same exterior coefficients. The positive-time
tail theorem shows why a bulk limit alone does not resolve the
positive-threshold problem.

\paragraph{Scope and evidence.}
This is an analytical synthesis of the operator foundation
\cite{OperatorProof,OperatorAssessment} and collision reconnaissance
\cite{OperatorCollision}. The former has an internal assessment within
its supplied-preparation scope; the latter has proof-author checks only.
The local collision statements are conditional on an event, except for
the explicitly constructed noncanonical control in Section~\ref{sec:control}.
That control satisfies the generic domain and heat identities and
nevertheless collides. The formulas therefore identify phenomena to
understand, rather than an independent canonical collision restriction.
All of them admit the same-information scalar Fourier/heat derivation.

We proceed from the complete source to the operator realization, then
from regular field evolution to local zero geometry and its quantitative
consequences. The critical coefficients are then related to the exact theta moments,
its half-line primitive and a contour representation of the exterior
factor. Interference, perturbation and higher multiplicity follow.
The real-zero benchmark is explicitly conditional; the exact control
shows why its curvature sign alone cannot exclude a collision.
Rodgers--Tao renormalization provides the connection to zero statistics.
The exact control and evidence boundaries lead to the outstanding tasks
in Section~\ref{sec:outstanding}.
